\section{Evaluation}\label{sec:eval}
\subsection{Setting}
Two sites are studied, both served by the one deployment described in Secs.~\ref{sec:arch}--\ref{sec:impl} from T0, 23~September~2026, 18:00~CEST; items from earlier pilots and from the acceptance checks before T0 are excluded. \textbf{S1} is the reference site, scaledaiops.org; its Respond stage is the scheduled agent, and every item, comment, label and pull request is public. \textbf{S2} is an anonymous B2B quick-commerce platform supplying restaurants, in its pre-launch phase, that embeds the same widget and routes items to its own private tracker, where its maintainers respond manually. It contributes the anonymised research rows of Sec.~\ref{sec:impl} only. Because capture, tracker and metric pipeline are identical on both sites, S2 is a comparison point for the Respond stage itself: agentic versus human first response.
\subsection{Protocol}
The protocol was fixed in the project's case-study document before data collection: a 12-week window from T0 (to 16~December~2026) with an interim read-out at week~4 (21~October~2026); per-kind, per-ISO-week metrics as defined in Sec.~\ref{sec:model}; targets stated in advance---TTFR under two hours (the agent's schedule), TTHR under 72~h at the median, TTC for bugs under 30 days at the median, closing notification for 100\,\% of consented closed items---and signal ratio and agent share reported without a target. The source of truth is, for S1, the public tracker read by the published \texttt{metrics}/\texttt{export} commands and, for S2, its weekly research rows; the resulting CSVs are committed with their export date. Items filed by the maintainer to test the system are labelled \texttt{spam} in the tracker and excluded.
\subsection{Results at week 4}
Table~\ref{tab:metrics} will report the interim metrics for both sites against the pre-registered targets; the values are filled from the committed week-4 export (\todo{2026-10-21}) and are deliberately blank until then. Operating cost (RQ1) will be reported from the actual cloud bill for the shared deployment; the agent runs within a subscription and has no per-item API charge.

\begin{table}[htbp]
\centering\footnotesize
\caption{Interim evaluation metrics (week 4). Targets are pre-registered; S2 responds without the agent.}
\label{tab:metrics}
\begin{tabular}{@{}l@{\hspace{0.75em}}l@{\hspace{0.75em}}r@{\hspace{0.75em}}r@{\hspace{0.75em}}r@{}}
\toprule
Metric & Target & S1 (bugs) & S1 (features) & S2 (all) \\
\midrule
$n$ (signal) & --- & \todo{} & \todo{} & \todo{} \\
TTFR (p50 / p90) & $<2$\,h & \todo{} & \todo{} & \todo{} \\
TTHR (p50) & $<72$\,h & \todo{} & \todo{} & \todo{} \\
TTC (p50) & $<30$\,d (bugs) & \todo{} & \todo{} & \todo{} \\
Loop closure & --- & \todo{} & \todo{} & \todo{} \\
Consented notification & 100\,\% & \todo{} & \todo{} & \todo{} \\
Signal ratio & --- & \todo{} & \todo{} & \todo{} \\
Agent share & --- & \todo{} & \todo{} & n/a \\
\bottomrule
\end{tabular}
\end{table}

\subsection{Ablation (RQ3)}
The dependency-minimal architecture is the result of removing components from an earlier design; we argue for the load-bearing role of what remains:
\begin{itemize}
    \item \textbf{Issue tracker vs.\ dedicated database:} the initial design had a relational store and an outbox. Removing them simplified the system and moved durability to the tracker; the trade-off is synchronous coupling---if the tracker is down, capture fails with an explicit \texttt{502} and the client retries safely with its idempotency key.
    \item \textbf{Sidecar vs.\ public personal data:} the private sidecar isolates e-mail addresses and screenshots; without it a public tracker could not be the system of record.
    \item \textbf{Bot check and honeypot:} protect the signal ratio by rejecting automated submissions before an issue is created; the counts of rejected and spam-labelled items are reported at week~4.
\end{itemize}
